% Options for packages loaded elsewhere
% Options for packages loaded elsewhere
\PassOptionsToPackage{unicode}{hyperref}
\PassOptionsToPackage{hyphens}{url}
\PassOptionsToPackage{dvipsnames,svgnames,x11names}{xcolor}
%
\documentclass[
  letterpaper,
  DIV=11,
  numbers=noendperiod]{scrartcl}
\usepackage{xcolor}
\usepackage{amsmath,amssymb}
\setcounter{secnumdepth}{-\maxdimen} % remove section numbering
\usepackage{iftex}
\ifPDFTeX
  \usepackage[T1]{fontenc}
  \usepackage[utf8]{inputenc}
  \usepackage{textcomp} % provide euro and other symbols
\else % if luatex or xetex
  \usepackage{unicode-math} % this also loads fontspec
  \defaultfontfeatures{Scale=MatchLowercase}
  \defaultfontfeatures[\rmfamily]{Ligatures=TeX,Scale=1}
\fi
\usepackage{lmodern}
\ifPDFTeX\else
  % xetex/luatex font selection
\fi
% Use upquote if available, for straight quotes in verbatim environments
\IfFileExists{upquote.sty}{\usepackage{upquote}}{}
\IfFileExists{microtype.sty}{% use microtype if available
  \usepackage[]{microtype}
  \UseMicrotypeSet[protrusion]{basicmath} % disable protrusion for tt fonts
}{}
\makeatletter
\@ifundefined{KOMAClassName}{% if non-KOMA class
  \IfFileExists{parskip.sty}{%
    \usepackage{parskip}
  }{% else
    \setlength{\parindent}{0pt}
    \setlength{\parskip}{6pt plus 2pt minus 1pt}}
}{% if KOMA class
  \KOMAoptions{parskip=half}}
\makeatother
% Make \paragraph and \subparagraph free-standing
\makeatletter
\ifx\paragraph\undefined\else
  \let\oldparagraph\paragraph
  \renewcommand{\paragraph}{
    \@ifstar
      \xxxParagraphStar
      \xxxParagraphNoStar
  }
  \newcommand{\xxxParagraphStar}[1]{\oldparagraph*{#1}\mbox{}}
  \newcommand{\xxxParagraphNoStar}[1]{\oldparagraph{#1}\mbox{}}
\fi
\ifx\subparagraph\undefined\else
  \let\oldsubparagraph\subparagraph
  \renewcommand{\subparagraph}{
    \@ifstar
      \xxxSubParagraphStar
      \xxxSubParagraphNoStar
  }
  \newcommand{\xxxSubParagraphStar}[1]{\oldsubparagraph*{#1}\mbox{}}
  \newcommand{\xxxSubParagraphNoStar}[1]{\oldsubparagraph{#1}\mbox{}}
\fi
\makeatother

\usepackage{color}
\usepackage{fancyvrb}
\newcommand{\VerbBar}{|}
\newcommand{\VERB}{\Verb[commandchars=\\\{\}]}
\DefineVerbatimEnvironment{Highlighting}{Verbatim}{commandchars=\\\{\}}
% Add ',fontsize=\small' for more characters per line
\usepackage{framed}
\definecolor{shadecolor}{RGB}{241,243,245}
\newenvironment{Shaded}{\begin{snugshade}}{\end{snugshade}}
\newcommand{\AlertTok}[1]{\textcolor[rgb]{0.68,0.00,0.00}{#1}}
\newcommand{\AnnotationTok}[1]{\textcolor[rgb]{0.37,0.37,0.37}{#1}}
\newcommand{\AttributeTok}[1]{\textcolor[rgb]{0.40,0.45,0.13}{#1}}
\newcommand{\BaseNTok}[1]{\textcolor[rgb]{0.68,0.00,0.00}{#1}}
\newcommand{\BuiltInTok}[1]{\textcolor[rgb]{0.00,0.23,0.31}{#1}}
\newcommand{\CharTok}[1]{\textcolor[rgb]{0.13,0.47,0.30}{#1}}
\newcommand{\CommentTok}[1]{\textcolor[rgb]{0.37,0.37,0.37}{#1}}
\newcommand{\CommentVarTok}[1]{\textcolor[rgb]{0.37,0.37,0.37}{\textit{#1}}}
\newcommand{\ConstantTok}[1]{\textcolor[rgb]{0.56,0.35,0.01}{#1}}
\newcommand{\ControlFlowTok}[1]{\textcolor[rgb]{0.00,0.23,0.31}{\textbf{#1}}}
\newcommand{\DataTypeTok}[1]{\textcolor[rgb]{0.68,0.00,0.00}{#1}}
\newcommand{\DecValTok}[1]{\textcolor[rgb]{0.68,0.00,0.00}{#1}}
\newcommand{\DocumentationTok}[1]{\textcolor[rgb]{0.37,0.37,0.37}{\textit{#1}}}
\newcommand{\ErrorTok}[1]{\textcolor[rgb]{0.68,0.00,0.00}{#1}}
\newcommand{\ExtensionTok}[1]{\textcolor[rgb]{0.00,0.23,0.31}{#1}}
\newcommand{\FloatTok}[1]{\textcolor[rgb]{0.68,0.00,0.00}{#1}}
\newcommand{\FunctionTok}[1]{\textcolor[rgb]{0.28,0.35,0.67}{#1}}
\newcommand{\ImportTok}[1]{\textcolor[rgb]{0.00,0.46,0.62}{#1}}
\newcommand{\InformationTok}[1]{\textcolor[rgb]{0.37,0.37,0.37}{#1}}
\newcommand{\KeywordTok}[1]{\textcolor[rgb]{0.00,0.23,0.31}{\textbf{#1}}}
\newcommand{\NormalTok}[1]{\textcolor[rgb]{0.00,0.23,0.31}{#1}}
\newcommand{\OperatorTok}[1]{\textcolor[rgb]{0.37,0.37,0.37}{#1}}
\newcommand{\OtherTok}[1]{\textcolor[rgb]{0.00,0.23,0.31}{#1}}
\newcommand{\PreprocessorTok}[1]{\textcolor[rgb]{0.68,0.00,0.00}{#1}}
\newcommand{\RegionMarkerTok}[1]{\textcolor[rgb]{0.00,0.23,0.31}{#1}}
\newcommand{\SpecialCharTok}[1]{\textcolor[rgb]{0.37,0.37,0.37}{#1}}
\newcommand{\SpecialStringTok}[1]{\textcolor[rgb]{0.13,0.47,0.30}{#1}}
\newcommand{\StringTok}[1]{\textcolor[rgb]{0.13,0.47,0.30}{#1}}
\newcommand{\VariableTok}[1]{\textcolor[rgb]{0.07,0.07,0.07}{#1}}
\newcommand{\VerbatimStringTok}[1]{\textcolor[rgb]{0.13,0.47,0.30}{#1}}
\newcommand{\WarningTok}[1]{\textcolor[rgb]{0.37,0.37,0.37}{\textit{#1}}}

\usepackage{longtable,booktabs,array}
\usepackage{calc} % for calculating minipage widths
% Correct order of tables after \paragraph or \subparagraph
\usepackage{etoolbox}
\makeatletter
\patchcmd\longtable{\par}{\if@noskipsec\mbox{}\fi\par}{}{}
\makeatother
% Allow footnotes in longtable head/foot
\IfFileExists{footnotehyper.sty}{\usepackage{footnotehyper}}{\usepackage{footnote}}
\makesavenoteenv{longtable}
\usepackage{graphicx}
\makeatletter
\newsavebox\pandoc@box
\newcommand*\pandocbounded[1]{% scales image to fit in text height/width
  \sbox\pandoc@box{#1}%
  \Gscale@div\@tempa{\textheight}{\dimexpr\ht\pandoc@box+\dp\pandoc@box\relax}%
  \Gscale@div\@tempb{\linewidth}{\wd\pandoc@box}%
  \ifdim\@tempb\p@<\@tempa\p@\let\@tempa\@tempb\fi% select the smaller of both
  \ifdim\@tempa\p@<\p@\scalebox{\@tempa}{\usebox\pandoc@box}%
  \else\usebox{\pandoc@box}%
  \fi%
}
% Set default figure placement to htbp
\def\fps@figure{htbp}
\makeatother





\setlength{\emergencystretch}{3em} % prevent overfull lines

\providecommand{\tightlist}{%
  \setlength{\itemsep}{0pt}\setlength{\parskip}{0pt}}



 


\KOMAoption{captions}{tableheading}
\makeatletter
\@ifpackageloaded{caption}{}{\usepackage{caption}}
\AtBeginDocument{%
\ifdefined\contentsname
  \renewcommand*\contentsname{Table of contents}
\else
  \newcommand\contentsname{Table of contents}
\fi
\ifdefined\listfigurename
  \renewcommand*\listfigurename{List of Figures}
\else
  \newcommand\listfigurename{List of Figures}
\fi
\ifdefined\listtablename
  \renewcommand*\listtablename{List of Tables}
\else
  \newcommand\listtablename{List of Tables}
\fi
\ifdefined\figurename
  \renewcommand*\figurename{Figure}
\else
  \newcommand\figurename{Figure}
\fi
\ifdefined\tablename
  \renewcommand*\tablename{Table}
\else
  \newcommand\tablename{Table}
\fi
}
\@ifpackageloaded{float}{}{\usepackage{float}}
\floatstyle{ruled}
\@ifundefined{c@chapter}{\newfloat{codelisting}{h}{lop}}{\newfloat{codelisting}{h}{lop}[chapter]}
\floatname{codelisting}{Listing}
\newcommand*\listoflistings{\listof{codelisting}{List of Listings}}
\makeatother
\makeatletter
\makeatother
\makeatletter
\@ifpackageloaded{caption}{}{\usepackage{caption}}
\@ifpackageloaded{subcaption}{}{\usepackage{subcaption}}
\makeatother
\usepackage{bookmark}
\IfFileExists{xurl.sty}{\usepackage{xurl}}{} % add URL line breaks if available
\urlstyle{same}
\hypersetup{
  pdftitle={National Malaria Elimination Program Dashboard: Technical Documentation and User Guide},
  pdfauthor={Enku Demssie},
  colorlinks=true,
  linkcolor={blue},
  filecolor={Maroon},
  citecolor={Blue},
  urlcolor={Blue},
  pdfcreator={LaTeX via pandoc}}


\title{National Malaria Elimination Program Dashboard: Technical
Documentation and User Guide}
\author{Enku Demssie}
\date{2026-06-01}
\begin{document}
\maketitle

\renewcommand*\contentsname{Table of contents}
{
\hypersetup{linkcolor=}
\setcounter{tocdepth}{4}
\tableofcontents
}

\section{Introduction}\label{introduction}

The National Malaria Elimination Program (NMEP) malaria dashboard was
developed to support routine monitoring and analysis of malaria
indicators using data from the national Health Management Information
System (HMIS), implemented through DHIS2.

The dashboard integrates routine HMIS data on malaria testing, confirmed
cases, species distribution, and malaria-related deaths across inpatient
and outpatient services. These data are used for monitoring malaria
trends over time and across geographic areas.

The system is designed to enable standardized analysis of key malaria
indicators at woreda level, supporting routine program monitoring,
reporting, and decision-making.

This document describes the step-by-step process used to develop the
dashboard, including data extraction, processing, indicator
construction, and deployment.

\section{Overview of The Dashboard Development
Process}\label{overview-of-the-dashboard-development-process}

The development of the HMIS-based NMEP dashboard followed a structured
workflow consisting of the following key steps:

\begin{enumerate}
\def\labelenumi{\arabic{enumi}.}
\tightlist
\item
  Data extraction from DHIS2
\item
  Data preparation and cleaning

  \begin{itemize}
  \tightlist
  \item
    Data preparation
  \item
    General cleaning and harmonization
  \item
    Geographic cleaning
  \end{itemize}
\item
  Dashboard dataset preparation

  \begin{itemize}
  \tightlist
  \item
    Historical data integration
  \item
    Lookup tables
  \item
    Time enrichment
  \item
    Aggregation
  \end{itemize}
\item
  Visualization and dashboard development
\item
  Deployment and routine updating
\end{enumerate}

Each step is described in detail in the following sections.

\section{1. Data Extraction from
DHIS2}\label{data-extraction-from-dhis2}

Data for the dashboard were extracted from the national Health
Management Information System (HMIS) implemented through DHIS2. The
extraction covered the period from January 2020 to April 2026 and was
conducted at the woreda level, which was selected as the lowest
administrative unit for analysis and visualization.

Data were drawn from three primary HMIS datasets to support
comprehensive analysis of malaria burden, service delivery, and
reporting performance.

\subsection{1.1 Disease surveilance
Data}\label{disease-surveilance-data}

The disease surveillance dataset was extracted from DHIS2 using the
following components:

\begin{itemize}
\tightlist
\item
  Data elements (malaria-related indicators)\\
\item
  Period\\
\item
  Organizational unit\\
\item
  Outcome (morbidity and mortality disaggregation)\\
\item
  Department (inpatient and outpatient disaggregation)
\end{itemize}

These components enabled structured disaggregation of malaria burden
across time, geography, and service delivery settings.

A total of 16 malaria-related data elements were included. These
comprised both ICD-10--based data elements (used prior to mid-2023) and
ICD-11--based data elements (currently in use). This approach ensured
continuity of analysis despite changes in disease classification
systems.

\textbf{Disease Surveillance Data Elements}

\textbf{ICD-10 (2018--2023)}

\begin{itemize}
\tightlist
\item
  B50--196 Malaria (Plasmodium falciparum malaria)\\
\item
  B50--197 Malaria (Plasmodium falciparum with cerebral complications)\\
\item
  B50--198 Malaria (Other severe and complicated Plasmodium falciparum
  malaria)\\
\item
  B50--199 Malaria (Plasmodium falciparum unspecified)\\
\item
  B50--200 Malaria (Plasmodium vivax malaria)\\
\item
  B50--201 Malaria (Plasmodium vivax with other complications)\\
\item
  B50--202 Malaria (Plasmodium vivax without complication)\\
\item
  B50--203 Mixed malaria (Other parasitologically confirmed malaria)\\
\item
  B50--204 Malaria (Unspecified malaria)
\end{itemize}

\textbf{ICD-11 (2023--present)}

\begin{itemize}
\item
  1F40 Malaria due to \emph{Plasmodium falciparum}\\
\item
  1F41 Malaria due to \emph{Plasmodium vivax}\\
\item
  1F42 Malaria due to \emph{Plasmodium malariae}\\
\item
  1F43 Malaria due to \emph{Plasmodium ovale}\\
\item
  1F44 Other parasitologically confirmed malaria\\
\item
  1F40/1F41 Mixed malaria (\emph{P. falciparum} and \emph{P. vivax})\\
\item
  1F45 Malaria without parasitological confirmation
\end{itemize}

\subsection{1.2 Service Delivery Data}\label{service-delivery-data}

Service delivery data were extracted separately from DHIS2, as this
dataset does not support the additional disaggregation dimensions used
in the disease surveillance dataset (i.e., outcome and department). As a
result, only the standard DHIS2 components were applied:

\begin{itemize}
\tightlist
\item
  Data elements\\
\item
  Period\\
\item
  Organizational unit
\end{itemize}

The dataset included three key malaria-related data elements:

\begin{itemize}
\tightlist
\item
  Malaria tests (legacy indicator used prior to 2023: ``Total number of
  slides or RDT performed for malaria diagnosis'')\\
\item
  Malaria tests (current indicator: ``Slides or RDT performed for
  malaria diagnosis'')\\
\item
  Confirmed positive malaria tests
\end{itemize}

This separation ensured consistency with the DHIS2 service reporting
structure while maintaining data integrity across indicator definitions
and system updates.

\subsection{1.3 Reporting Completeness and Timeliness
Data}\label{reporting-completeness-and-timeliness-data}

Data on service report completeness and timeliness were extracted from
the integrated service reporting dataset used across malaria, NTD, and
NCD programs, as no malaria-specific reporting performance dataset is
available.

The dataset included nine data elements, disaggregated by facility type:

\begin{itemize}
\tightlist
\item
  Hospitals, health centers, and private facilities\\
\item
  Basic health posts\\
\item
  Comprehensive health posts
\end{itemize}

For each facility type, the following indicators were extracted:

\begin{itemize}
\tightlist
\item
  Actual reports submitted\\
\item
  Reports submitted on time\\
\item
  Expected reports
\end{itemize}

These data elements were used to assess service report completeness and
timeliness across different levels of the health system.

\textbf{Data Export Format}

All datasets were manually exported from DHIS2 in CSV format. These
exports formed the input datasets for subsequent data cleaning,
transformation, and indicator calculation within the dashboard
development workflow.

\section{2. Data Preparation \&
cleaning}\label{data-preparation-cleaning}

The raw HMIS datasets (disease surveillance, service delivery, and
reporting completeness \& timeliness) were processed to create a unified
analysis-ready dataset for the dashboard. The cleaning and
transformation process was implemented in R.

\subsection{2.1 Data preparation}\label{data-preparation}

\subsubsection{2.1.1 Standardization of variable names and geographic
fields}\label{standardization-of-variable-names-and-geographic-fields}

The three datasets extracted from DHIS2 were first read into the R
environment.

\begin{Shaded}
\begin{Highlighting}[]
\FunctionTok{library}\NormalTok{(tidyverse) }
\FunctionTok{library}\NormalTok{(lubridate)}

\CommentTok{\# reading the disease report dat}
\NormalTok{df\_species }\OtherTok{\textless{}{-}} \FunctionTok{suppressWarnings}\NormalTok{(}
  \FunctionTok{read\_csv}\NormalTok{(}\StringTok{"data/raw/April 2026 hmis/disease\_dat\_april2026.csv"}\NormalTok{, }\AttributeTok{show\_col\_types =} \ConstantTok{FALSE}\NormalTok{))}

\CommentTok{\# reading the service report data}
\NormalTok{df\_service }\OtherTok{\textless{}{-}} \FunctionTok{suppressWarnings}\NormalTok{(}\FunctionTok{read\_csv}\NormalTok{(}\StringTok{"data/raw/April 2026 hmis/service\_dat\_april2026.csv"}\NormalTok{, }\AttributeTok{show\_col\_types =} \ConstantTok{FALSE}\NormalTok{))}

\CommentTok{\# reading the dataset for service report completeness \& timeliness}
\NormalTok{df\_report }\OtherTok{\textless{}{-}}\FunctionTok{suppressMessages}\NormalTok{(}
  \FunctionTok{suppressWarnings}\NormalTok{(}
  \FunctionTok{read\_csv}\NormalTok{(}\StringTok{"data/raw/April 2026 hmis/report\_dat\_april2026.csv"}\NormalTok{,}\AttributeTok{show\_col\_types =} \ConstantTok{FALSE}\NormalTok{)))}

\CommentTok{\# checking the variable names for the 3 dataset}
\FunctionTok{names}\NormalTok{(df\_species)}
\FunctionTok{names}\NormalTok{(df\_service)}
\FunctionTok{names}\NormalTok{(df\_report)}

\CommentTok{\# preparing a function to rename the adminstartive levels }
\NormalTok{rename\_geo }\OtherTok{\textless{}{-}} \ControlFlowTok{function}\NormalTok{(df)\{}
\NormalTok{  df }\SpecialCharTok{|\textgreater{}}
    \FunctionTok{rename}\NormalTok{(}
      \AttributeTok{period =}\NormalTok{ periodname,}
      \AttributeTok{region =}\NormalTok{ orgunitlevel2,}
      \AttributeTok{zone   =}\NormalTok{ orgunitlevel3,}
      \AttributeTok{woreda =}\NormalTok{ orgunitlevel4}
\NormalTok{    )}
\NormalTok{\}}

\CommentTok{\# performing demonstrative leve renaming using the above function}
\NormalTok{df\_species }\OtherTok{\textless{}{-}} \FunctionTok{rename\_geo}\NormalTok{(df\_species)}
\NormalTok{df\_service }\OtherTok{\textless{}{-}} \FunctionTok{rename\_geo}\NormalTok{(df\_service)}
\NormalTok{df\_report  }\OtherTok{\textless{}{-}} \FunctionTok{rename\_geo}\NormalTok{(df\_report)}
\end{Highlighting}
\end{Shaded}

At this stage, variable names and administrative unit fields were
standardized across all datasets to ensure consistency prior to
integration.

\subsubsection{2.1.2 Reshaping to long
format}\label{reshaping-to-long-format}

Each dataset was transformed from wide format to long format to enable
indicator-level analysis.

\begin{Shaded}
\begin{Highlighting}[]
\NormalTok{df\_species\_long }\OtherTok{\textless{}{-}}\NormalTok{ df\_species }\SpecialCharTok{|\textgreater{}}
  \FunctionTok{pivot\_longer}\NormalTok{(}\AttributeTok{cols =} \FunctionTok{matches}\NormalTok{(}\StringTok{"ESV|B50"}\NormalTok{),}
               \AttributeTok{names\_to =} \StringTok{"data\_type"}\NormalTok{,}
               \AttributeTok{values\_to =} \StringTok{"value"}\NormalTok{)}

\NormalTok{df\_service\_long }\OtherTok{\textless{}{-}}\NormalTok{ df\_service }\SpecialCharTok{|\textgreater{}}
  \FunctionTok{pivot\_longer}\NormalTok{(}\AttributeTok{cols =} \FunctionTok{contains}\NormalTok{(}\StringTok{"RDT"}\NormalTok{),}
               \AttributeTok{names\_to =} \StringTok{"data\_type"}\NormalTok{,}
               \AttributeTok{values\_to =} \StringTok{"value"}\NormalTok{)}

\NormalTok{df\_report\_long }\OtherTok{\textless{}{-}}\NormalTok{ df\_report }\SpecialCharTok{|\textgreater{}}
  \FunctionTok{pivot\_longer}\NormalTok{(}\AttributeTok{cols =} \FunctionTok{matches}\NormalTok{(}\StringTok{"NCD|NTD"}\NormalTok{),}
               \AttributeTok{names\_to =} \StringTok{"data\_type"}\NormalTok{,}
               \AttributeTok{values\_to =} \StringTok{"value"}\NormalTok{)}
\end{Highlighting}
\end{Shaded}

\subsubsection{2.1.3 Standardization of indicator
dimensions}\label{standardization-of-indicator-dimensions}

To allow integration, missing categorical variables were added to ensure
a consistent structure across datasets.

\begin{Shaded}
\begin{Highlighting}[]
\NormalTok{df\_service\_long }\OtherTok{\textless{}{-}}\NormalTok{ df\_service\_long }\SpecialCharTok{|\textgreater{}}
  \FunctionTok{mutate}\NormalTok{(}\AttributeTok{department =} \ConstantTok{NA\_character\_}\NormalTok{,}
         \AttributeTok{outcome =} \ConstantTok{NA\_character\_}\NormalTok{,}
         \AttributeTok{facility\_type =} \ConstantTok{NA\_character\_}\NormalTok{)}

\NormalTok{df\_report\_long }\OtherTok{\textless{}{-}}\NormalTok{ df\_report\_long }\SpecialCharTok{|\textgreater{}}
  \FunctionTok{mutate}\NormalTok{(}\AttributeTok{department =} \ConstantTok{NA\_character\_}\NormalTok{,}
         \AttributeTok{outcome =} \ConstantTok{NA\_character\_}\NormalTok{)}

\NormalTok{df\_species\_long }\OtherTok{\textless{}{-}}\NormalTok{ df\_species\_long }\SpecialCharTok{|\textgreater{}}
  \FunctionTok{mutate}\NormalTok{(}\AttributeTok{facility\_type =} \ConstantTok{NA\_character\_}\NormalTok{)}
\end{Highlighting}
\end{Shaded}

\subsubsection{2.1.4 Dataset integration}\label{dataset-integration}

After standardizing variable names and aligning dataset structures, the
disease surveillance, service delivery, and reporting datasets were
integrated into a single unified dataset.

\begin{Shaded}
\begin{Highlighting}[]
\NormalTok{combined\_hmis }\OtherTok{\textless{}{-}} \FunctionTok{bind\_rows}\NormalTok{(}
\NormalTok{  df\_species\_long,}
\NormalTok{  df\_service\_long,}
\NormalTok{  df\_report\_long}
\NormalTok{)}
\end{Highlighting}
\end{Shaded}

The resulting dataset (`combined\_hmis`) serves as the primary
integrated dataset for subsequent processing steps, including geographic
cleaning and indicator construction for the dashboard.

\begin{Shaded}
\begin{Highlighting}[]
\CommentTok{\# saving the integrated dataset}
\FunctionTok{saveRDS}\NormalTok{(combined\_hmis,}
        \StringTok{"data/raw/April 2026 hmis/combined{-}hmis{-}apr\_2026.rds"}\NormalTok{)}
\end{Highlighting}
\end{Shaded}

\subsection{2.2. General Data Cleaning and
Harmonization}\label{general-data-cleaning-and-harmonization}

Following the creation of the integrated dataset
(\texttt{combined\_hmis}), a general cleaning and harmonization process
was applied to standardize variable names, data elements, time
variables, and geographic identifiers prior to geographic cleaning.

This step ensured that the dataset was consistent, structured, and ready
for geospatial standardization and downstream analysis.

\subsubsection{2.2.1 Loading processed
dataset}\label{loading-processed-dataset}

The integrated HMIS dataset from the previous step was loaded into the R
environment for further cleaning and transformation.

\begin{Shaded}
\begin{Highlighting}[]
\CommentTok{\# libraries }
\FunctionTok{library}\NormalTok{(tidyverse) }
\FunctionTok{library}\NormalTok{(lubridate) }
\FunctionTok{library}\NormalTok{(janitor) }
\FunctionTok{library}\NormalTok{(sf) }
\FunctionTok{library}\NormalTok{(stringr)}

\CommentTok{\#reading the integrated dataset}
\NormalTok{combined\_hmis }\OtherTok{\textless{}{-}} \FunctionTok{readRDS}\NormalTok{( }\StringTok{"data/raw/April 2026 hmis/combined{-}hmis{-}apr\_2026.rds"}\NormalTok{ )}
\end{Highlighting}
\end{Shaded}

\subsubsection{2.2.2 Standardization of variable names and text
fields}\label{standardization-of-variable-names-and-text-fields}

Variable names and geographic text fields were cleaned to ensure
consistency across administrative units.

\begin{Shaded}
\begin{Highlighting}[]
\NormalTok{combined\_hmis }\OtherTok{\textless{}{-}}\NormalTok{ combined\_hmis }\SpecialCharTok{|\textgreater{}}
\NormalTok{  janitor}\SpecialCharTok{::}\FunctionTok{clean\_names}\NormalTok{() }\SpecialCharTok{|\textgreater{}}
  \FunctionTok{mutate}\NormalTok{(}
    \AttributeTok{region =} \FunctionTok{gsub}\NormalTok{(}\StringTok{"Region"}\NormalTok{, }\StringTok{""}\NormalTok{, region),}
    \AttributeTok{region =} \FunctionTok{gsub}\NormalTok{(}\StringTok{"region"}\NormalTok{, }\StringTok{""}\NormalTok{, region),}
    \AttributeTok{region =} \FunctionTok{gsub}\NormalTok{(}\StringTok{"City Administration"}\NormalTok{, }\StringTok{""}\NormalTok{, region),}
    \AttributeTok{region =} \FunctionTok{gsub}\NormalTok{(}\StringTok{"al Health Bureau"}\NormalTok{, }\StringTok{""}\NormalTok{, region),}
    \AttributeTok{region =} \FunctionTok{gsub}\NormalTok{(}\StringTok{"Ethiopian"}\NormalTok{, }\StringTok{""}\NormalTok{, region),}
    \AttributeTok{region =} \FunctionTok{gsub}\NormalTok{(}\StringTok{"Ethiopia"}\NormalTok{, }\StringTok{""}\NormalTok{, region),}
    \AttributeTok{region =} \FunctionTok{gsub}\NormalTok{(}\StringTok{"  "}\NormalTok{, }\StringTok{" "}\NormalTok{, region),}
    \AttributeTok{region =} \FunctionTok{str\_trim}\NormalTok{(region)}
\NormalTok{  ) }\SpecialCharTok{|\textgreater{}}
  \FunctionTok{mutate}\NormalTok{(}
    \AttributeTok{zone =} \FunctionTok{gsub}\NormalTok{(}\StringTok{"Health Office"}\NormalTok{, }\StringTok{""}\NormalTok{, zone),}
    \AttributeTok{zone =} \FunctionTok{gsub}\NormalTok{(}\StringTok{"Zone"}\NormalTok{, }\StringTok{""}\NormalTok{, zone),}
    \AttributeTok{zone =} \FunctionTok{gsub}\NormalTok{(}\StringTok{"City Administration"}\NormalTok{, }\StringTok{""}\NormalTok{, zone),}
    \AttributeTok{zone =} \FunctionTok{gsub}\NormalTok{(}\StringTok{"Zone Health Department"}\NormalTok{, }\StringTok{""}\NormalTok{, zone),}
    \AttributeTok{zone =} \FunctionTok{gsub}\NormalTok{(}\StringTok{"Health Department"}\NormalTok{, }\StringTok{""}\NormalTok{, zone),}
    \AttributeTok{zone =} \FunctionTok{gsub}\NormalTok{(}\StringTok{"Health department"}\NormalTok{, }\StringTok{""}\NormalTok{, zone),}
    \AttributeTok{zone =} \FunctionTok{gsub}\NormalTok{(}\StringTok{"Woreda"}\NormalTok{, }\StringTok{""}\NormalTok{, zone),}
    \AttributeTok{zone =} \FunctionTok{gsub}\NormalTok{(}\StringTok{"Zonal"}\NormalTok{, }\StringTok{""}\NormalTok{, zone),}
    \AttributeTok{woreda =} \FunctionTok{gsub}\NormalTok{(}\StringTok{"Town Adminstration"}\NormalTok{, }\StringTok{""}\NormalTok{, woreda),}
    \AttributeTok{woreda =} \FunctionTok{gsub}\NormalTok{(}\StringTok{"Health Office"}\NormalTok{, }\StringTok{""}\NormalTok{, woreda),}
    \AttributeTok{woreda =} \FunctionTok{str\_squish}\NormalTok{(woreda)}
\NormalTok{  )}
\end{Highlighting}
\end{Shaded}

\subsubsection{2.2.3 Harmonization of data
elements}\label{harmonization-of-data-elements}

Malaria-related indicators and reporting variables were standardized
into a simplified classification system to enable consistent analysis
across ICD-10 and ICD-11 periods.

\begin{Shaded}
\begin{Highlighting}[]
\NormalTok{pf\_conf }\OtherTok{\textless{}{-}} \FunctionTok{c}\NormalTok{(}\StringTok{"B50{-}196 Malaria (Plasmodium falciparum malaria)"}\NormalTok{,}
             \StringTok{"B50{-}197 Malaria (Plasmodium falciparum malaria with cerebral complications)"}\NormalTok{,}
             \StringTok{"B50{-}198 Malaria (Other severe and complicated Plasmodium falciparum malaria)"}\NormalTok{,}
             \StringTok{"B50{-}199 Malaria (Plasmodium falciparum malaria unspecified)"}\NormalTok{,}
             \StringTok{"ESV{-}ICD11 1F40 {-} Malaria due to Plasmodium falciparum"}\NormalTok{,}
             \StringTok{"ESV{-}ICD11 1F44 {-} Other parasitologically confirmed malaria"}\NormalTok{,}
             \StringTok{"B50{-}203 Mixed Malaria (Other parasitologically confirmed malaria)"}
\NormalTok{)}

\NormalTok{pv\_conf }\OtherTok{\textless{}{-}} \FunctionTok{c}\NormalTok{(}\StringTok{"B50{-}200 Malaria (Plasmodium vivax malaria)"}\NormalTok{,}
             \StringTok{"B50{-}201 Malaria (Plasmodium vivax malaria with other complications)"}\NormalTok{,}
             \StringTok{"B50{-}202 Malaria (Plasmodium vivax malaria without complication)"}\NormalTok{,}
             \StringTok{"ESV{-}ICD11 1F41 {-} Malaria due to Plasmodium vivax"}
\NormalTok{)}

\NormalTok{mixed\_conf }\OtherTok{\textless{}{-}} \FunctionTok{c}\NormalTok{(}\StringTok{"ESV{-}ICD11 1F40/1F41 {-}  Malaria due to Plasmodium falciparum associated with Malaria due to Plasmodium Vivax (Mixed Malaria)"}\NormalTok{,}
                \StringTok{"B50{-}203 Mixed Malaria (Other parasitologically confirmed malaria)"}
\NormalTok{)}

\NormalTok{pm\_conf }\OtherTok{\textless{}{-}} \FunctionTok{c}\NormalTok{(}\StringTok{"ESV{-}ICD11 1F42 {-} Malaria due to Plasmodium malariae"}\NormalTok{)}

\NormalTok{po\_conf }\OtherTok{\textless{}{-}} \FunctionTok{c}\NormalTok{(}\StringTok{"ESV{-}ICD11 1F43 {-} Malaria due to Plasmodium ovale"}\NormalTok{)}

\NormalTok{clinical\_malaria }\OtherTok{\textless{}{-}} \FunctionTok{c}\NormalTok{(}\StringTok{"B50{-}204 Malaria (Unspecified malaria)"}\NormalTok{,}
                      \StringTok{"ESV{-}ICD11 1F45 {-} Malaria without parasitological confirmation"}\NormalTok{)}

\NormalTok{tested }\OtherTok{\textless{}{-}} \FunctionTok{c}\NormalTok{(}\StringTok{"MAL\_Slides or RDT performed for malaria diagnosis"}\NormalTok{,}
            \StringTok{"{-}2017\_Total number of slides or RDT performed for malaria diagnosis"}\NormalTok{)}

\NormalTok{positives }\OtherTok{\textless{}{-}} \FunctionTok{c}\NormalTok{(}\StringTok{"MAL\_Slides or RDT Positive"}\NormalTok{)}
\NormalTok{actual\_reports }\OtherTok{\textless{}{-}} \FunctionTok{c}\NormalTok{(}\StringTok{"08.1 {-} Malaria, NTD, NCD | Basic Health post | Monthly {-} Actual reports"}\NormalTok{,}
                    \StringTok{"08.2 {-} Malaria, NTD, NCD | Hospital, Health center, Clinic | Monthly {-} Actual reports"}\NormalTok{,}
                    \StringTok{"08.3 {-} Malaria, NTD, NCD | Comprehensive HP | Monthly {-} Actual reports"}\NormalTok{)}

\NormalTok{actual\_on\_time }\OtherTok{\textless{}{-}} \FunctionTok{c}\NormalTok{(}\StringTok{"08.1 {-} Malaria, NTD, NCD | Basic Health post | Monthly {-} Actual reports on time"}\NormalTok{,}
                    \StringTok{"08.2 {-} Malaria, NTD, NCD | Hospital, Health center, Clinic | Monthly {-} Actual reports on time"}\NormalTok{,}
                    \StringTok{"08.3 {-} Malaria, NTD, NCD | Comprehensive HP | Monthly {-} Actual reports on time"}\NormalTok{)}

\NormalTok{expected }\OtherTok{\textless{}{-}} \FunctionTok{c}\NormalTok{(}\StringTok{"08.1 {-} Malaria, NTD, NCD | Basic Health post | Monthly {-} Expected reports"}\NormalTok{,}
              \StringTok{"08.2 {-} Malaria, NTD, NCD | Hospital, Health center, Clinic | Monthly {-} Expected reports"}\NormalTok{,}
              \StringTok{"08.3 {-} Malaria, NTD, NCD | Comprehensive HP | Monthly {-} Expected reports"}\NormalTok{)}
\end{Highlighting}
\end{Shaded}

\subsubsection{2.2.4 Standardization of data
elements}\label{standardization-of-data-elements}

All indicators were standardized into simplified categories for
analysis.

\begin{Shaded}
\begin{Highlighting}[]
\NormalTok{combined\_hmis }\OtherTok{\textless{}{-}}\NormalTok{ combined\_hmis }\SpecialCharTok{|\textgreater{}}
  \FunctionTok{mutate}\NormalTok{(}
    \AttributeTok{data\_type =} \FunctionTok{case\_when}\NormalTok{(}
\NormalTok{      data\_type }\SpecialCharTok{\%in\%}\NormalTok{ pf\_conf }\SpecialCharTok{\textasciitilde{}} \StringTok{"pf\_conf"}\NormalTok{,}
\NormalTok{      data\_type }\SpecialCharTok{\%in\%}\NormalTok{ pv\_conf }\SpecialCharTok{\textasciitilde{}} \StringTok{"pv\_conf"}\NormalTok{,}
\NormalTok{      data\_type }\SpecialCharTok{\%in\%}\NormalTok{ positives }\SpecialCharTok{\textasciitilde{}} \StringTok{"positives"}\NormalTok{,}
\NormalTok{      data\_type }\SpecialCharTok{\%in\%}\NormalTok{ tested }\SpecialCharTok{\textasciitilde{}} \StringTok{"tested"}\NormalTok{,}
\NormalTok{      data\_type }\SpecialCharTok{\%in\%}\NormalTok{ clinical\_malaria }\SpecialCharTok{\textasciitilde{}} \StringTok{"clinical"}\NormalTok{,}
\NormalTok{      data\_type }\SpecialCharTok{\%in\%}\NormalTok{ mixed\_conf }\SpecialCharTok{\textasciitilde{}} \StringTok{"mixed\_conf"}\NormalTok{,}
\NormalTok{      data\_type }\SpecialCharTok{\%in\%}\NormalTok{ po\_conf }\SpecialCharTok{\textasciitilde{}} \StringTok{"po\_conf"}\NormalTok{,}
\NormalTok{      data\_type }\SpecialCharTok{\%in\%}\NormalTok{ pm\_conf }\SpecialCharTok{\textasciitilde{}} \StringTok{"pm\_conf"}\NormalTok{,}
\NormalTok{      data\_type }\SpecialCharTok{\%in\%}\NormalTok{ actual\_reports }\SpecialCharTok{\textasciitilde{}} \StringTok{"actual\_reports"}\NormalTok{,}
\NormalTok{      data\_type }\SpecialCharTok{\%in\%}\NormalTok{ actual\_on\_time }\SpecialCharTok{\textasciitilde{}} \StringTok{"actual\_reports\_ontime"}\NormalTok{,}
\NormalTok{      data\_type }\SpecialCharTok{\%in\%}\NormalTok{ expected }\SpecialCharTok{\textasciitilde{}} \StringTok{"expected\_reports"}\NormalTok{,}
      \ConstantTok{TRUE} \SpecialCharTok{\textasciitilde{}}\NormalTok{ data\_type}
\NormalTok{    )}
\NormalTok{  )}
\end{Highlighting}
\end{Shaded}

\subsubsection{2.2.5 Temporal standardization (Ethiopian to Gregorian
calendar)}\label{temporal-standardization-ethiopian-to-gregorian-calendar}

The period variable was parsed to extract Ethiopian calendar components
and to derive corresponding Gregorian calendar dates. This was done to
enable flexible time-series analysis while retaining the original
Ethiopian calendar format for reference and reporting consistency.

Both calendar systems were retained in the final dataset to support
different analytical and reporting needs.

\begin{Shaded}
\begin{Highlighting}[]
\NormalTok{eth\_months }\OtherTok{\textless{}{-}} \FunctionTok{c}\NormalTok{(}\StringTok{"Meskerem"}\NormalTok{,}\StringTok{"Tikemet"}\NormalTok{,}\StringTok{"Hidar"}\NormalTok{,}\StringTok{"Tahesas"}\NormalTok{,}
                \StringTok{"Tir"}\NormalTok{,}\StringTok{"Yekatit"}\NormalTok{,}\StringTok{"Megabit"}\NormalTok{,}\StringTok{"Miazia"}\NormalTok{,}
                \StringTok{"Ginbot"}\NormalTok{,}\StringTok{"Sene"}\NormalTok{,}\StringTok{"Hamle"}\NormalTok{,}\StringTok{"Nehase"}\NormalTok{)}

\NormalTok{greg\_months }\OtherTok{\textless{}{-}} \FunctionTok{c}\NormalTok{(}\DecValTok{9}\NormalTok{,}\DecValTok{10}\NormalTok{,}\DecValTok{11}\NormalTok{,}\DecValTok{12}\NormalTok{,}\DecValTok{1}\NormalTok{,}\DecValTok{2}\NormalTok{,}\DecValTok{3}\NormalTok{,}\DecValTok{4}\NormalTok{,}\DecValTok{5}\NormalTok{,}\DecValTok{6}\NormalTok{,}\DecValTok{7}\NormalTok{,}\DecValTok{8}\NormalTok{)}
\NormalTok{month\_map }\OtherTok{\textless{}{-}} \FunctionTok{setNames}\NormalTok{(}\DecValTok{1}\SpecialCharTok{:}\DecValTok{12}\NormalTok{, eth\_months)}

\NormalTok{combined\_hmis }\OtherTok{\textless{}{-}}\NormalTok{ combined\_hmis }\SpecialCharTok{\%\textgreater{}\%}
  \FunctionTok{mutate}\NormalTok{(}
    \AttributeTok{eth\_month =} \FunctionTok{str\_extract}\NormalTok{(period, }\StringTok{"\^{}[\^{} ]+"}\NormalTok{),}
    \AttributeTok{eth\_year  =} \FunctionTok{as.integer}\NormalTok{(}\FunctionTok{str\_extract}\NormalTok{(period, }\StringTok{"}\SpecialCharTok{\textbackslash{}\textbackslash{}}\StringTok{d+$"}\NormalTok{)),}
    \AttributeTok{eth\_month\_num =}\NormalTok{ month\_map[eth\_month],}
    \AttributeTok{greg\_year =} \FunctionTok{if\_else}\NormalTok{(eth\_month\_num }\SpecialCharTok{\textless{}=} \DecValTok{4}\NormalTok{, eth\_year }\SpecialCharTok{+} \DecValTok{7}\NormalTok{, eth\_year }\SpecialCharTok{+} \DecValTok{8}\NormalTok{),}
    \AttributeTok{greg\_month =}\NormalTok{ greg\_months[eth\_month\_num],}
    \AttributeTok{greg\_date =} \FunctionTok{make\_date}\NormalTok{(greg\_year, greg\_month, }\DecValTok{1}\NormalTok{)}
\NormalTok{  ) }\SpecialCharTok{|\textgreater{}}
  \FunctionTok{mutate}\NormalTok{(}
    \AttributeTok{greg\_month =} \FunctionTok{month}\NormalTok{(greg\_date, }\AttributeTok{label =} \ConstantTok{TRUE}\NormalTok{),}
    \AttributeTok{greg\_year =} \FunctionTok{year}\NormalTok{(greg\_date),}
    \AttributeTok{month\_year =} \FunctionTok{paste}\NormalTok{(greg\_month, }\StringTok{"{-}"}\NormalTok{, greg\_year)}
\NormalTok{  ) }\SpecialCharTok{|\textgreater{}}
  \FunctionTok{select}\NormalTok{(}
\NormalTok{    region, zone, woreda,}
\NormalTok{    eth\_month, eth\_year,}
\NormalTok{    greg\_date, greg\_month, greg\_year, month\_year,}
\NormalTok{    facility\_type, department, outcome, data\_type, value}
\NormalTok{  )}
\end{Highlighting}
\end{Shaded}

The dataset was saved after general cleaning and harmonization for use
in subsequent steps, including geographic cleaning and downstream
analysis.

\begin{Shaded}
\begin{Highlighting}[]
\FunctionTok{saveRDS}\NormalTok{(}
\NormalTok{  combined\_hmis,}
  \StringTok{"data/processed/april 2026/apr\_2026\_hmis\_general\_cleaning\_completed.rds"}
\NormalTok{)}
\end{Highlighting}
\end{Shaded}

\subsection{2.3 Geographic name
Cleaning}\label{geographic-name-cleaning}

Following general data cleaning and harmonization, geographic variables
were further standardized to ensure alignment between HMIS
administrative unit names and the official Ethiopian administrative
boundary reference.

Geographic cleaning was performed separately for each administrative
level (region, zone, and woreda) using a standardized cleaning function
developed by the PATH/MACEPA data analytics team. This function was used
to harmonize geographic identifiers across datasets and align them with
the reference administrative structure.

This step was implemented prior to indicator generation and
visualization and was applied sequentially across all three
administrative levels.

\textbf{Management of unmatched geographic names}

The geographic cleaning functions compare HMIS administrative unit names
against the official administrative boundary reference dataset and the
corresponding preprocessed correction files. In cases where unmatched
names are identified, the unmatched records are exported for manual
review.

These records are then assessed and mapped to the appropriate
administrative unit names. The resulting corrections are added to the
relevant preprocessing file (e.g., region, zone, or woreda correction
file).

Once the correction file has been updated, the cleaning function can be
rerun to apply the new mappings automatically. This process creates a
growing repository of validated name corrections, reducing manual effort
in future dashboard updates and improving consistency across reporting
periods.

\subsubsection{2.3.1 Region-level geographic
cleaning}\label{region-level-geographic-cleaning}

Region-level geographic cleaning was performed on the dataset resulting
from general cleaning and harmonization. The objective was to
standardize region names to match the official Ethiopian administrative
reference.

This step ensured that all region identifiers in the HMIS dataset were
consistent with the official administrative boundary definitions derived
from the national shapefile reference.

\emph{Inputs to the region cleaning process}

Three inputs were used in the region-level cleaning procedure:

\begin{itemize}
\item
  \textbf{HMIS dataset (input data):}\\
  The output from general cleaning and harmonization\\
  (\texttt{apr\_2026\_hmis\_general\_cleaning\_completed.rds})
\item
  \textbf{Preprocessed region name corrections:}\\
  A manually prepared mapping file containing inconsistent region names
  and their corresponding corrected values\\
  (\texttt{region\_names\_corrected.csv})
\item
  \textbf{Reference region names:}\\
  official administrative boundary reference dataset\\
  (\texttt{eth-admin3-v1082-clean-id.rds})
\end{itemize}

\textbf{Region cleaning process}

Region cleaning was performed using a standardized function developed by
the PATH/MACEPA data analytics team. This function applies predefined
matching logic, correction mappings, and reference validation to align
region names with the official administrative structure.

\begin{Shaded}
\begin{Highlighting}[]
\CommentTok{\# libraries}
\FunctionTok{library}\NormalTok{(tidyverse)}
\FunctionTok{library}\NormalTok{(lubridate)}
\FunctionTok{library}\NormalTok{(janitor)}

\CommentTok{\# sourcing the geographic cleaning function}
\FunctionTok{source}\NormalTok{(}\StringTok{"functions/geog\_cleaning\_accessories.R"}\NormalTok{)}

\CommentTok{\# paths for the files needed }
\NormalTok{ref\_path }\OtherTok{\textless{}{-}} \StringTok{"data/eth{-}admin3{-}v1082{-}clean{-}id.rds"}
\NormalTok{dat\_path }\OtherTok{\textless{}{-}} \StringTok{"data/processed/april 2026/apr\_2026\_hmis\_general\_cleaning\_completed.rds"}
\NormalTok{pre\_path }\OtherTok{\textless{}{-}} \StringTok{"data/preprocessed/region\_names\_corrected.csv"}
\NormalTok{out\_path }\OtherTok{\textless{}{-}} \StringTok{"data/processed/april 2026/apr\_2026\_hmis\_region\_cleaning\_completed.rds"}

\CommentTok{\# applying region cleaning function}
\NormalTok{rclean }\OtherTok{\textless{}{-}} \FunctionTok{region\_clean}\NormalTok{(}
  \AttributeTok{dat\_path =}\NormalTok{ dat\_path,}
  \AttributeTok{pre\_path =}\NormalTok{ pre\_path,}
  \AttributeTok{ref\_path =}\NormalTok{ ref\_path,}
  \AttributeTok{out\_path =}\NormalTok{ out\_path}
\NormalTok{)}
\end{Highlighting}
\end{Shaded}

\begin{itemize}
\tightlist
\item
  \textbf{Output of region cleaning}
\end{itemize}

The output of this step is a region-standardized HMIS dataset in which
all region names are aligned with the official administrative reference.

This dataset serves as the input for subsequent zone-level geographic
cleaning.

\subsubsection{2.3.2 Zone-level geographic
cleaning}\label{zone-level-geographic-cleaning}

Following region name standardization, zone names were cleaned and
standardized to ensure alignment with the official administrative
boundary reference dataset used in the dashboard.

This step was performed after region cleaning because zone matching
relies on a consistent regional structure.

\emph{Inputs to the zone cleaning process}

Three inputs were used in the zone-level cleaning procedure:

\begin{itemize}
\item
  \textbf{HMIS dataset (input data):}\\
  Output from the region cleaning step\\
  (\texttt{apr\_2026\_hmis\_region\_cleaning\_completed.rds})
\item
  \textbf{Preprocessed zone name corrections:}\\
  A manually prepared mapping file containing inconsistent zone names
  and their corresponding corrected values\\
  (\texttt{zone\_names\_corrected.csv})
\item
  \textbf{Reference administrative boundary dataset:}\\
  Official administrative names used as the reference for geographic
  standardization\\
  (\texttt{eth-admin3-v1082-clean-id.rds})
\end{itemize}

\textbf{Zone cleaning process}

Zone cleaning was performed using the standardized function developed by
the PATH/MACEPA data analytics team. The function applies predefined
correction mappings and reference matching procedures to align zone
names with the official administrative hierarchy.

\begin{Shaded}
\begin{Highlighting}[]
\CommentTok{\# libraries}
\FunctionTok{library}\NormalTok{(tidyverse)}
\FunctionTok{library}\NormalTok{(lubridate)}
\FunctionTok{library}\NormalTok{(janitor)}

\CommentTok{\# sourcing the geographic cleaning function}
\FunctionTok{source}\NormalTok{(}\StringTok{"functions/geog\_cleaning\_accessories.R"}\NormalTok{)}

\CommentTok{\# path for the reference administrative boundary dataset}
\NormalTok{ref\_path }\OtherTok{\textless{}{-}} \StringTok{"data/eth{-}admin3{-}v1082{-}clean{-}id.rds"}

\CommentTok{\# path for the input dataset}
\NormalTok{dat\_path }\OtherTok{\textless{}{-}} \StringTok{"data/processed/april 2026/apr\_2026\_hmis\_region\_cleaning\_completed.rds"}

\CommentTok{\# path for preprocessed zone corrections}
\NormalTok{pre\_path }\OtherTok{\textless{}{-}} \StringTok{"data/preprocessed/zone\_names\_corrected.csv"}

\CommentTok{\# path for the output dataset}
\NormalTok{out\_path }\OtherTok{\textless{}{-}} \StringTok{"data/processed/april 2026/apr\_2026\_hmis\_zone\_cleaning\_completed.rds"}

\CommentTok{\# clean zone names}
\NormalTok{zclean }\OtherTok{\textless{}{-}} \FunctionTok{zone\_clean}\NormalTok{(}
\NormalTok{  dat\_path,}
\NormalTok{  pre\_path,}
\NormalTok{  ref\_path,}
\NormalTok{  out\_path}
\NormalTok{)}
\end{Highlighting}
\end{Shaded}

\subsubsection{2.3.3 Woreda-level geographic
cleaning}\label{woreda-level-geographic-cleaning}

Following zone-level standardization, woreda names were cleaned and
standardized to ensure alignment with the official administrative
boundary reference.

This step was performed after region and zone cleaning because woreda
matching relies on a consistent administrative hierarchy. Standardizing
woreda names ensures that HMIS records can be accurately linked to the
corresponding administrative boundaries used for analysis and
visualization.

\emph{Inputs to the woreda cleaning process}

Three inputs were used in the woreda-level cleaning procedure:

\begin{itemize}
\item
  \textbf{HMIS dataset (input data):}\\
  Output from the zone cleaning step\\
  (\texttt{apr\_2026\_hmis\_zone\_cleaning\_completed.rds})
\item
  \textbf{Preprocessed woreda name corrections:}\\
  A cumulative correction file containing previously validated mappings
  between unmatched HMIS woreda names and the corresponding official
  administrative names\\
  (\texttt{woreda\_names\_corrected.csv})
\item
  \textbf{Reference administrative boundary dataset:}\\
  Official administrative names used as the reference for geographic
  standardization\\
  (\texttt{eth-admin3-v1082-clean-id.rds})
\end{itemize}

\textbf{Woreda cleaning process}

Woreda cleaning was performed using the same standardized function
developed by the PATH/MACEPA team. The function applies predefined
correction mappings and reference matching procedures to align woreda
names with the official administrative hierarchy.

\begin{Shaded}
\begin{Highlighting}[]
\CommentTok{\# libraries }
\FunctionTok{library}\NormalTok{(tidyverse) }
\FunctionTok{library}\NormalTok{(lubridate) }
\FunctionTok{library}\NormalTok{(janitor)}

\CommentTok{\# sourcing the geographic cleaning function}
\FunctionTok{source}\NormalTok{(}\StringTok{"functions/geog\_cleaning\_accessories.R"}\NormalTok{)}

\CommentTok{\# path for the reference administrative boundary dataset}
\NormalTok{ref\_path }\OtherTok{\textless{}{-}} \StringTok{"data/eth{-}admin3{-}v1082{-}clean{-}id.rds"}

\CommentTok{\# path for the input dataset}
\NormalTok{dat\_path }\OtherTok{\textless{}{-}} \StringTok{"data/processed/april 2026/apr\_2026\_hmis\_zone\_cleaning\_completed.rds"}

\CommentTok{\# path for preprocessed woreda corrections}
\NormalTok{pre\_path }\OtherTok{\textless{}{-}} \StringTok{"data/preprocessed/woreda\_names\_corrected.csv"}

\CommentTok{\# path for the output dataset}
\NormalTok{out\_path }\OtherTok{\textless{}{-}} \StringTok{"data/processed/april 2026/apr\_2026\_hmis\_woreda\_cleaning\_completed.rds"}

\CommentTok{\# clean woreda names}
\NormalTok{wclean }\OtherTok{\textless{}{-}} \FunctionTok{woreda\_clean}\NormalTok{( dat\_path, pre\_path, ref\_path, out\_path )}
\end{Highlighting}
\end{Shaded}

\textbf{Integration of administrative identifiers}

After woreda names were standardized, the cleaned HMIS dataset was
linked to the administrative boundary reference dataset to incorporate
the unique woreda identifier (\texttt{id\_1082}).

The \texttt{id\_1082} variable is a standardized administrative code
assigned to each woreda in the national administrative boundary dataset.
Incorporating this identifier provides a stable geographic reference
that is not affected by spelling variations or naming inconsistencies.

The identifier is used throughout the dashboard workflow for geographic
aggregation, linkage with auxiliary datasets, and spatial visualization.

\begin{Shaded}
\begin{Highlighting}[]
\CommentTok{\# read cleaned HMIS data}
\NormalTok{clean\_woredas }\OtherTok{\textless{}{-}} \FunctionTok{readRDS}\NormalTok{(out\_path) }\SpecialCharTok{|\textgreater{}}
  \FunctionTok{select}\NormalTok{(}
    \SpecialCharTok{{-}}\NormalTok{region\_old,}
    \SpecialCharTok{{-}}\NormalTok{zone\_old,}
    \SpecialCharTok{{-}}\NormalTok{woreda\_old,}
    \SpecialCharTok{{-}}\NormalTok{non\_zero\_ever,}
    \SpecialCharTok{{-}}\NormalTok{matched}
\NormalTok{  )}

\CommentTok{\# read reference administrative boundary dataset}
\NormalTok{sf }\OtherTok{\textless{}{-}} \FunctionTok{readRDS}\NormalTok{(ref\_path) }\SpecialCharTok{|\textgreater{}}
\NormalTok{  sf}\SpecialCharTok{::}\FunctionTok{st\_drop\_geometry}\NormalTok{() }\SpecialCharTok{|\textgreater{}}
  \FunctionTok{select}\NormalTok{(}\SpecialCharTok{{-}}\NormalTok{region\_old, }\SpecialCharTok{{-}}\NormalTok{pop\_2022)}

\CommentTok{\# join administrative identifier}
\NormalTok{final\_out }\OtherTok{\textless{}{-}}\NormalTok{ clean\_woredas }\SpecialCharTok{|\textgreater{}}
  \FunctionTok{left\_join}\NormalTok{(}
\NormalTok{    sf,}
    \AttributeTok{by =} \FunctionTok{c}\NormalTok{(}\StringTok{"region"}\NormalTok{, }\StringTok{"zone"}\NormalTok{, }\StringTok{"woreda"}\NormalTok{)}
\NormalTok{  )}

\CommentTok{\# save final geographic{-}cleaned dataset}
\FunctionTok{saveRDS}\NormalTok{(final\_out, out\_path)}
\end{Highlighting}
\end{Shaded}

\textbf{Output of woreda cleaning}

The output of this step is a dataset with standardized region, zone, and
woreda names aligned with the official administrative boundary reference
dataset.

This dataset serves as the geographic-cleaned HMIS dataset and provides
the basis for subsequent aggregation, indicator calculation,
visualization, and dashboard development.

\section{3. dashboard-ready dataset
preparation}\label{dashboard-ready-dataset-preparation}

\subsection{3.1 Integration with the historical HMIS
dataset}\label{integration-with-the-historical-hmis-dataset}

The dashboard is designed to display malaria trends over multiple years.
Therefore, each newly processed monthly HMIS dataset must be integrated
with the existing historical HMIS dataset maintained for the dashboard.

At this stage, the processed HMIS dataset for the current reporting
period was appended to the historical HMIS dataset. This process creates
an updated master dataset containing all available HMIS records required
for dashboard analysis and visualization.

\begin{Shaded}
\begin{Highlighting}[]
\CommentTok{\# libraries}
\FunctionTok{library}\NormalTok{(tidyverse)}
\FunctionTok{library}\NormalTok{(lubridate)}

\CommentTok{\# read historical HMIS dataset}
\NormalTok{hmis\_large }\OtherTok{\textless{}{-}} \FunctionTok{readRDS}\NormalTok{(}
  \StringTok{"data/processed/jan 2020{-}mar 2026/clean\_hmis\_2020\_2026.rds"}
\NormalTok{)}

\CommentTok{\# read newly processed HMIS dataset}
\NormalTok{hmis\_apr\_2026 }\OtherTok{\textless{}{-}} \FunctionTok{readRDS}\NormalTok{(}
  \StringTok{"data/processed/april 2026/clean\_apr\_2026.rds"}
\NormalTok{)}

\CommentTok{\# append new data to historical dataset}
\NormalTok{binded\_hmis }\OtherTok{\textless{}{-}} \FunctionTok{bind\_rows}\NormalTok{(}
\NormalTok{  hmis\_large,}
\NormalTok{  hmis\_apr\_2026}
\NormalTok{)}

\CommentTok{\# save updated master dataset}
\FunctionTok{saveRDS}\NormalTok{(}
\NormalTok{  binded\_hmis,}
  \StringTok{"data/processed/april 2026/upto\_apr\_2026\_hmis\_for\_app.rds"}
\NormalTok{)}
\end{Highlighting}
\end{Shaded}

\textbf{Output of dataset integration}

The output of this step is an updated master HMIS dataset containing
both the historical records and the newly processed reporting period.

This dataset serves as the source dataset for subsequent aggregation and
indicator calculation steps used in the dashboard.

\subsection{3.2 Creation of lookup tables and analytical
dataset}\label{creation-of-lookup-tables-and-analytical-dataset}

Following geographic cleaning, a set of supporting lookup tables were
prepared to facilitate dashboard functionality and indicator
calculation.

The dashboard uses two types of lookup tables:

\begin{itemize}
\item
  \textbf{Static lookup tables}, which are created once and reused
  across routine dashboard updates. These include the population lookup
  and time lookup tables.
\item
  \textbf{Dynamic lookup tables}, which are generated from the latest
  HMIS data and therefore require updating whenever a new HMIS
  extraction is processed. This includes the geographic lookup table.
\end{itemize}

These lookup tables are subsequently used to enrich the cleaned HMIS
dataset and create the final analysis-ready dataset used throughout the
dashboard.

\textbf{Geographic lookup table}

A geographic lookup table was generated from the geographically cleaned
HMIS dataset. The lookup contains unique administrative identifiers and
names for each region, zone, and woreda represented in the data.

Unlike the population and time lookup tables, the geographic lookup is
derived directly from the HMIS dataset and should be regenerated
whenever a new HMIS extraction is processed.

\begin{Shaded}
\begin{Highlighting}[]
\CommentTok{\#libraries}
\FunctionTok{library}\NormalTok{(tidyverse)}
\FunctionTok{library}\NormalTok{(lubridate)}
\FunctionTok{library}\NormalTok{(sf)}

\CommentTok{\# reading the cleaned hmis data}
\NormalTok{hmis }\OtherTok{\textless{}{-}} \FunctionTok{readRDS}\NormalTok{(}\StringTok{"data/processed/april 2026/april\_mar\_2026\_hmis\_woreda\_cleaning\_completed.rds"}\NormalTok{)}

\CommentTok{\# read sf}
\NormalTok{sf }\OtherTok{\textless{}{-}} \FunctionTok{readRDS}\NormalTok{(}\StringTok{"data/eth{-}admin3{-}v1082{-}clean{-}id.rds"}\NormalTok{)}\SpecialCharTok{|\textgreater{}}
\NormalTok{  sf}\SpecialCharTok{::}\FunctionTok{st\_drop\_geometry}\NormalTok{()}


\CommentTok{\# prepare a geo look up}
\NormalTok{geo\_lookup }\OtherTok{\textless{}{-}}\NormalTok{ hmis }\SpecialCharTok{\%\textgreater{}\%}
  \FunctionTok{distinct}\NormalTok{(}
\NormalTok{    id\_1082,}
\NormalTok{    region,}
\NormalTok{    zone,}
\NormalTok{    woreda}
\NormalTok{  ) }\SpecialCharTok{\%\textgreater{}\%}
  \FunctionTok{arrange}\NormalTok{(region, zone, woreda)}

\CommentTok{\# save the geo look up}
\FunctionTok{saveRDS}\NormalTok{(geo\_lookup, }\StringTok{"data/processed/april 2026/geo\_lookup.rds"}\NormalTok{)}
\end{Highlighting}
\end{Shaded}

\textbf{Population lookup table}

A population lookup table containing annual woreda-level population
estimates was prepared for use in population-based indicator
calculations.

This lookup is maintained separately from the HMIS data and can be
reused across routine dashboard updates.

\begin{Shaded}
\begin{Highlighting}[]
\CommentTok{\# create a population lookup}
\NormalTok{pop\_lookup }\OtherTok{\textless{}{-}} \FunctionTok{read\_csv}\NormalTok{(}\StringTok{"data/eth{-}admin3{-}v1082{-}clean{-}id\_with\_pop\_2030.csv"}\NormalTok{)}\SpecialCharTok{|\textgreater{}}
  \FunctionTok{filter}\NormalTok{(year }\SpecialCharTok{\%in\%} \FunctionTok{c}\NormalTok{(}\DecValTok{2020}\SpecialCharTok{:}\DecValTok{2026}\NormalTok{))}

\CommentTok{\# save population lookup (saved once, no need to repeat)}
\CommentTok{\#saveRDS(pop\_lookup, "data/processed/population\_lookup.rds")}
\end{Highlighting}
\end{Shaded}

\textbf{Time lookup table}

A time lookup table was created to provide standardized Gregorian and
Ethiopian Fiscal Year (EFY) date variables used throughout the
dashboard.

This lookup is created independently of the HMIS data and can be reused
across routine dashboard updates. It only requires modification if the
dashboard period is extended beyond the range covered by the existing
lookup table.

\begin{Shaded}
\begin{Highlighting}[]
\CommentTok{\# creating time lookup}
\NormalTok{time\_lookup }\OtherTok{\textless{}{-}} \FunctionTok{tibble}\NormalTok{(}
  \AttributeTok{greg\_date =} \FunctionTok{seq}\NormalTok{(}
    \AttributeTok{from =} \FunctionTok{as.Date}\NormalTok{(}\StringTok{"2020{-}01{-}01"}\NormalTok{),}
    \AttributeTok{to =} \FunctionTok{as.Date}\NormalTok{(}\StringTok{"2030{-}12{-}01"}\NormalTok{),}
    \AttributeTok{by =} \StringTok{"month"}
\NormalTok{  )}
\NormalTok{)}

\CommentTok{\# create gregorian variable}
\NormalTok{time\_lookup }\OtherTok{\textless{}{-}}\NormalTok{ time\_lookup }\SpecialCharTok{\%\textgreater{}\%}
  \FunctionTok{mutate}\NormalTok{(}
    \AttributeTok{greg\_year =} \FunctionTok{year}\NormalTok{(greg\_date),}
    
    \AttributeTok{greg\_month =} \FunctionTok{month}\NormalTok{(}
\NormalTok{      greg\_date,}
      \AttributeTok{label =} \ConstantTok{TRUE}\NormalTok{,}
      \AttributeTok{abbr =} \ConstantTok{TRUE}
\NormalTok{    ),}
    \CommentTok{\# character type gregorian year for dashboard UI}
   \AttributeTok{greg\_year\_ui =} \FunctionTok{factor}\NormalTok{(}
\NormalTok{        greg\_year,}
        \AttributeTok{levels =} \DecValTok{2020}\SpecialCharTok{:}\DecValTok{2026}
\NormalTok{      ),}
    
    \AttributeTok{greg\_month\_num =} \FunctionTok{month}\NormalTok{(greg\_date),}
    
    \AttributeTok{month\_year =} \FunctionTok{format}\NormalTok{(}
\NormalTok{      greg\_date,}
      \StringTok{"\%b{-}\%Y"}
\NormalTok{    )}
\NormalTok{  )}

\CommentTok{\# create EFY variable}
\NormalTok{time\_lookup }\OtherTok{\textless{}{-}}\NormalTok{ time\_lookup }\SpecialCharTok{\%\textgreater{}\%}
  \FunctionTok{mutate}\NormalTok{(}
    
    \AttributeTok{efy =} \FunctionTok{ifelse}\NormalTok{(}
\NormalTok{      greg\_month\_num }\SpecialCharTok{\textgreater{}=} \DecValTok{7}\NormalTok{,}
\NormalTok{      greg\_year }\SpecialCharTok{{-}} \DecValTok{7}\NormalTok{,}
\NormalTok{      greg\_year }\SpecialCharTok{{-}} \DecValTok{8}
\NormalTok{    ),}
    
    \AttributeTok{efy\_month\_num =} \FunctionTok{case\_when}\NormalTok{(}
\NormalTok{      greg\_month\_num }\SpecialCharTok{\textgreater{}=} \DecValTok{7} \SpecialCharTok{\textasciitilde{}}\NormalTok{ greg\_month\_num }\SpecialCharTok{{-}} \DecValTok{6}\NormalTok{,}
      \ConstantTok{TRUE} \SpecialCharTok{\textasciitilde{}}\NormalTok{ greg\_month\_num }\SpecialCharTok{+} \DecValTok{6}
\NormalTok{    )}
\NormalTok{  )}


\CommentTok{\# creating EFY display label}
\NormalTok{time\_lookup }\OtherTok{\textless{}{-}}\NormalTok{ time\_lookup }\SpecialCharTok{\%\textgreater{}\%}
  \FunctionTok{mutate}\NormalTok{(}
    \AttributeTok{efy\_time =}\NormalTok{ month\_year}
\NormalTok{  )}

\CommentTok{\# convert to factor}
\NormalTok{time\_lookup }\OtherTok{\textless{}{-}}\NormalTok{ time\_lookup }\SpecialCharTok{\%\textgreater{}\%}
  \FunctionTok{mutate}\NormalTok{(}
    
    \AttributeTok{greg\_month =} \FunctionTok{factor}\NormalTok{(}
\NormalTok{      greg\_month,}
      \AttributeTok{levels =} \FunctionTok{c}\NormalTok{(}
        \StringTok{"Jan"}\NormalTok{,}\StringTok{"Feb"}\NormalTok{,}\StringTok{"Mar"}\NormalTok{,}\StringTok{"Apr"}\NormalTok{,}
        \StringTok{"May"}\NormalTok{,}\StringTok{"Jun"}\NormalTok{,}\StringTok{"Jul"}\NormalTok{,}\StringTok{"Aug"}\NormalTok{,}
        \StringTok{"Sep"}\NormalTok{,}\StringTok{"Oct"}\NormalTok{,}\StringTok{"Nov"}\NormalTok{,}\StringTok{"Dec"}
\NormalTok{      )}
\NormalTok{    ),}
    
    \AttributeTok{month\_year =} \FunctionTok{factor}\NormalTok{(}
\NormalTok{      month\_year,}
      \AttributeTok{levels =} \FunctionTok{format}\NormalTok{(}
\NormalTok{        greg\_date,}
        \StringTok{"\%b{-}\%Y"}
\NormalTok{      )}
\NormalTok{    )}
\NormalTok{  )}

\CommentTok{\# save the time lookup (saved once, no need to repeat)}
\CommentTok{\#saveRDS(time\_lookup,"data/processed/time\_lookup.rds")}
\end{Highlighting}
\end{Shaded}

\subsection{3.3 Temporal enrichment of the HMIS
dataset}\label{temporal-enrichment-of-the-hmis-dataset}

After creating the time lookup table, it was joined to the
geographically cleaned and integrated HMIS dataset using the Gregorian
date variable (\texttt{greg\_date}).

This step enriches the HMIS data with additional time-related variables
required for dashboard filtering, aggregation, and visualization. These
include Gregorian year and month variables, Ethiopian Fiscal Year (EFY)
variables, and user-interface-friendly date fields.

By storing these variables in the HMIS dataset, time transformations do
not need to be recalculated repeatedly during dashboard execution,
improving efficiency and ensuring consistent date handling throughout
the application.

\begin{Shaded}
\begin{Highlighting}[]
\CommentTok{\# join time lookup to the geography cleaned hmis}
\NormalTok{hmis }\OtherTok{\textless{}{-}}\NormalTok{ hmis }\SpecialCharTok{\%\textgreater{}\%}
  \FunctionTok{left\_join}\NormalTok{(}
\NormalTok{    time\_lookup,}
    \AttributeTok{by =} \StringTok{"greg\_date"}
\NormalTok{  )}

\CommentTok{\# selecting the columns needed for subsequent processing}
\NormalTok{hmis }\OtherTok{\textless{}{-}}\NormalTok{ hmis }\SpecialCharTok{|\textgreater{}}
  \FunctionTok{select}\NormalTok{(}
\NormalTok{    id\_1082,}
\NormalTok{    region,}
\NormalTok{    zone,}
\NormalTok{    woreda,}
\NormalTok{    greg\_date,}
    \AttributeTok{greg\_year =}\NormalTok{ greg\_year.y,}
    \AttributeTok{greg\_month =}\NormalTok{ greg\_month.y,}
\NormalTok{    greg\_month\_num,}
    \AttributeTok{month\_year =}\NormalTok{ month\_year.y,}
\NormalTok{    efy,}
\NormalTok{    greg\_year\_ui,}
\NormalTok{    efy\_month\_num,}
\NormalTok{    efy\_time,}
\NormalTok{    facility\_type,}
\NormalTok{    department,}
\NormalTok{    outcome,}
\NormalTok{    data\_type,}
\NormalTok{    value}
\NormalTok{  )}
\end{Highlighting}
\end{Shaded}

\textbf{Output of lookup table creation and temporal enrichment}

The output of this step is a geographically standardized and temporally
enriched HMIS dataset.

The dataset contains harmonized malaria indicators together with
standardized administrative identifiers and Gregorian/Ethiopian Fiscal
Year (EFY) time variables required for subsequent processing.

This dataset serves as the input for the aggregation step.

\subsection{3.4 Aggregation}\label{aggregation}

To improve dashboard performance and reduce computational load in the
Shiny application, the integrated HMIS dataset was aggregated prior to
deployment.

This step reduces data redundancy by summarizing records at the lowest
analytical level required for visualization while preserving all key
analytical dimensions.

\begin{Shaded}
\begin{Highlighting}[]
\CommentTok{\# libraries}
\FunctionTok{library}\NormalTok{(tidyverse)}
\FunctionTok{library}\NormalTok{(lubridate)}
\FunctionTok{library}\NormalTok{(sf)}

\CommentTok{\# read integrated HMIS dataset}
\NormalTok{hmis }\OtherTok{\textless{}{-}} \FunctionTok{readRDS}\NormalTok{(}
  \StringTok{"data/processed/april 2026/upto\_apr\_2026\_hmis\_for\_app.rds"}
\NormalTok{)}

\CommentTok{\# aggregate HMIS data for dashboard efficiency}
\NormalTok{hmis\_dashboard }\OtherTok{\textless{}{-}}\NormalTok{ hmis }\SpecialCharTok{|\textgreater{}}
  \FunctionTok{group\_by}\NormalTok{(}
\NormalTok{    id\_1082,}
\NormalTok{    region,}
\NormalTok{    zone,}
\NormalTok{    woreda,}
\NormalTok{    facility\_type,}
\NormalTok{    greg\_date,}
\NormalTok{    department,}
\NormalTok{    outcome,}
\NormalTok{    data\_type}
\NormalTok{  ) }\SpecialCharTok{|\textgreater{}}
  \FunctionTok{summarise}\NormalTok{(}
    \AttributeTok{value =} \FunctionTok{sum}\NormalTok{(value, }\AttributeTok{na.rm =} \ConstantTok{TRUE}\NormalTok{),}
    \AttributeTok{.groups =} \StringTok{"drop"}
\NormalTok{  )}
\end{Highlighting}
\end{Shaded}

\textbf{Output of aggregation}

The output of this step is a reduced and performance-optimized dataset
(hmis\_dashboard) in which duplicate records have been aggregated while
preserving all analytical dimensions required for dashboard
visualizations.

This improves dashboard responsiveness by reducing memory usage and
computational overhead during filtering and rendering.

\section{4. Visualization and Dashboard
Development}\label{visualization-and-dashboard-development}

\subsection{4.1 Application initialization and data
loading}\label{application-initialization-and-data-loading}

Before launching the user interface and server components, the
application loads all required datasets, lookup tables, geographic
boundary files, and supporting objects.

This initialization step ensures that all data required for filtering,
indicator calculation, and visualization are available when the
dashboard starts.

The following resources are loaded:

\begin{itemize}
\tightlist
\item
  Dashboard-ready HMIS dataset
\item
  Population lookup table
\item
  Geographic lookup table
\item
  Administrative boundary shapefile
\end{itemize}

In addition, geographic layers are prepared for woreda-, zone-, and
region-level mapping, and application-wide settings such as date
ordering and visual themes are initialized.

\subsection{4.2 Dashboard structure and
navigation}\label{dashboard-structure-and-navigation}

The dashboard was developed as a web-based interactive application using
R Shiny. The application is organized into two main pages designed to
support both user orientation and data exploration.

\textbf{Overview page}

The Overview page serves as the entry point to the dashboard and
provides background information to help users understand the purpose,
scope, and interpretation of the dashboard outputs.

The page includes the following sections:

\begin{itemize}
\tightlist
\item
  About the dashboard
\item
  Data sources
\item
  Key indicators
\item
  Methodological notes and
\item
  Dashboard user guide
\end{itemize}

This page is intended to provide contextual information and support
consistent interpretation of dashboard outputs.

\textbf{Dashboard page}

The Dashboard page contains all interactive visualizations and
analytical outputs. It is organized into a sidebar panel containing
user-defined filters and a main panel displaying a map, key performance
indicators (KPIs), and charts.

\textbf{Sidebar filters}

The sidebar contains three groups of filters:

\emph{Geographic filters}

\begin{itemize}
\tightlist
\item
  Region
\item
  Zone
\item
  Woreda
\end{itemize}

These filters allow users to view data at different administrative
levels and geographic locations.

\emph{Time filters}

\begin{itemize}
\tightlist
\item
  Gregorian year
\item
  Ethiopian Fiscal Year (EFY)
\item
  Month
\end{itemize}

These filters enable temporal analysis of malaria indicators.

\emph{Reporting filters}

\begin{itemize}
\tightlist
\item
  Facility type
\end{itemize}

The facility type filter is used only for reporting completeness and
timeliness indicators. All other filters apply across all dashboard
components.

\textbf{Dashboard outputs}

The main dashboard panel is organized into four rows of analytical
outputs.

\emph{Row 1}

\begin{itemize}
\tightlist
\item
  Geographic map
\item
  Key Performance Indicators (KPIs)
\end{itemize}

\emph{Row 2}

\begin{itemize}
\tightlist
\item
  Malaria testing and Test Positivity Rate (TPR)
\item
  Species proportion
\end{itemize}

The testing and TPR visualization combines malaria testing volume as a
bar chart with TPR displayed as a line on a secondary axis.

\emph{Row 3}

\begin{itemize}
\tightlist
\item
  Parasite incidence per 1,000 population
\item
  Malaria mortality rate per 100,000 population at risk
\end{itemize}

\emph{Row 4}

\begin{itemize}
\tightlist
\item
  Reporting completeness
\item
  Reporting timeliness
\end{itemize}

\textbf{Visualization modes}

All dashboard charts include a display toggle that allows users to
switch between:

\begin{itemize}
\tightlist
\item
  Aggregate view and
\item
  Monthly trend view
\end{itemize}

The aggregate view summarizes indicator values over the selected period
and geography, while the monthly view displays temporal trends across
individual months.

\section{5. Indicator Definitions and
Calculations}\label{indicator-definitions-and-calculations}

The dashboard presents key malaria surveillance, testing, mortality,
species distribution, and reporting performance indicators derived from
HMIS data. Indicator definitions and calculation methods are described
below.

\subsection{5.1 Malaria Cases}\label{malaria-cases}

\textbf{Confirmed Malaria Cases}

Confirmed malaria cases include all laboratory-confirmed malaria
infections reported through HMIS, regardless of species.

\textbf{Clinical Malaria Cases}

Clinical malaria cases refer to malaria diagnoses made without
parasitological confirmation.

\textbf{All Malaria Cases}

All malaria cases represent the total malaria burden reported through
HMIS and are calculated as:

\emph{All Cases = Confirmed Cases + Clinical Cases}

\subsection{5.2 Malaria Deaths}\label{malaria-deaths}

Malaria deaths are defined as inpatient malaria-related deaths reported
through the HMIS disease surveillance system.

\subsection{5.3 Test Positivity Rate
(TPR)}\label{test-positivity-rate-tpr}

The Test Positivity Rate (TPR) measures the proportion of malaria tests
that yield a positive result.

\emph{TPR = Positive Malaria Tests/Total malaria test performed x 100}

Where:\\
- Positive malaria tests = laboratory-confirmed positive malaria tests\\
- Total malaria tests performed = all malaria diagnostic tests
(microscopy and/or RDT)

\subsection{5.4 Species Distribution}\label{species-distribution}

Species proportions describe the relative contribution of each malaria
species among confirmed malaria cases.

\textbf{Plasmodium falciparum Proportion} = Pf cases/Total confirmed
cases*100

\textbf{Plasmodium vivax Proportion} = Pv cases/Total confirmed
cases*100

\textbf{Mixed malaria (pf/pv) Proportion} = Mixed cases/Total confirmed
cases*100

\emph{Total Confirmed Cases = Pf + Pv + Mixed + P. malariae + P. ovale}

\subsection{5.5 Annual Parasite Incidence
(API)}\label{annual-parasite-incidence-api}

The dashboard uses adjusted malaria cases to account for clinically
diagnosed malaria cases.

Adjusted malaria cases are calculated as:

\emph{Adjusted Cases = Confirmed Cases + (TPR × Clinical Cases)}

Annual Parasite Incidence (API) is then calculated as:

\emph{API = Adjusted Cases/ Total Population × 1000} and expressed as
cases per 1,000 population.

\subsection{5.6 Annualized Monthly Parasite
Incidence}\label{annualized-monthly-parasite-incidence}

For monthly dashboard displays, incidence is annualized to facilitate
comparison with annual incidence targets.

\emph{Annualized monthly incidence = (Adjusted cases X n/12)/Total
population X 1000}

Where:\\
\emph{n = annualization factor used to project the observed reporting
period to a 12-month equivalent}

\subsection{5.7 Malaria Mortality Rate}\label{malaria-mortality-rate}

The malaria mortality rate is calculated using the population at risk as
the denominator.

In accordance with the National Malaria Strategic Plan,the population at
risk is defined as:~ \emph{Population at Risk = Total Population × 69\%}

The mortality rate is calculated as:

\emph{Motality\_rate = In-patient malaria deaths/population at risk X
100,000} and expressed as deaths per 100,000 population at risk.

\subsection{5.8 Annualized Monthly Mortality
Rate}\label{annualized-monthly-mortality-rate}

For monthly dashboard displays, mortality is annualized using the same
approach applied to incidence.

\emph{Annualized monthly mortality = (inpatient malaria deaths X
n/12)/population at risk X 100,000}

Where:\\
\emph{n = annualization factor used to project the observed reporting
period to a 12-month equivalent}

\subsection{5.9 Reporting Completeness}\label{reporting-completeness}

Reporting completeness measures the proportion of expected reports that
were submitted.

\emph{Reporting completeness = Actual reports/Expected reports X 100}

\subsection{5.10 Reporting Timeliness}\label{reporting-timeliness}

Reporting timeliness measures the proportion of expected reports that
were submitted on time.

\emph{Reporting timeliness = Reports received on time/Expected reports X
100}

\textbf{Notes on Reporting Indicators}

Reporting completeness and reporting timeliness are assessed using only
the \textbf{HMIS service reporting dataset}. Both the numerator and
denominator values are system-generated by DHIS2 and are not calculated
manually within the dashboard.

These indicators assess reporting performance rather than malaria burden
and are displayed separately from surveillance and service delivery
indicators.

\section{6. Monthly Dashboard Update
Procedure}\label{monthly-dashboard-update-procedure}

The dashboard is designed to be updated monthly as new HMIS data become
available. The update process involves downloading the latest HMIS data,
running the existing data processing scripts, replacing the dashboard
dataset, and republishing the application.

Near-complete HMIS data are typically available in DHIS2 from
approximately the 22nd day of the following month. Dashboard updates
should therefore be performed after this date to ensure adequate
reporting completeness.

\textbf{The monthly update workflow consists of the following steps:}

\begin{enumerate}
\def\labelenumi{\arabic{enumi}.}
\tightlist
\item
  Download the latest HMIS datasets from DHIS2:

  \begin{itemize}
  \item
    Disease surveillance dataset
  \item
    Service delivery dataset
  \item
    Reporting completeness and timeliness dataset

    Save the downloaded files in the designated raw data folder.
  \end{itemize}
\item
  Run the data processing scripts sequentially:

  \begin{itemize}
  \item
    Data preparation script
  \item
    General data cleaning and harmonization script
  \item
    Region-level geographic cleaning script
  \item
    Zone-level geographic cleaning script
  \item
    Woreda-level geographic cleaning script
  \item
    Dashboard dataset preparation script
  \item
    Historical data integration script
  \item
    Dashboard aggregation script
  \end{itemize}
\item
  Review any unmatched geographic names generated during the cleaning
  process.

  \begin{itemize}
  \item
    Update the corresponding correction files if new unmatched names are
    identified:

    \begin{itemize}
    \item
      region\_names\_corrected.csv
    \item
      zone\_names\_corrected.csv
    \item
      woreda\_names\_corrected.csv
    \end{itemize}
  \end{itemize}
\item
  Rerun the relevant geographic cleaning script(s) if correction files
  are updated.
\item
  Generate the updated dashboard-ready dataset.
\item
  Replace the existing dashboard dataset with the newly generated
  dataset.
\item
  Run the Shiny application and verify that:

  \begin{itemize}
  \item
    The latest reporting month is available.
  \item
    Dashboard visualizations render correctly.
  \item
    Key indicators and maps display expected values.
  \item
    Filters function as intended.
  \end{itemize}
\item
  Republish the dashboard to ShinyApps.io. Perform a final review of the
  published dashboard to confirm that the updated data are visible and
  functioning correctly.
\end{enumerate}

This workflow should be repeated each month to maintain an up-to-date
dashboard for routine malaria surveillance and program monitoring.




\end{document}
